\subsection*{61. The promised key result and a first list of axioms (pp.151--156)}
Grothendieck resumes the comparison of a functor $i:A\to(Cat)$ with its induced $i^*:(Cat)\to\widehat A$, temporarily allowing two saturated arrow classes $\underline W$ and $\underline W'$. The desired equivalence of localizations is strengthened to the exact inverse-image condition $\underline W'=f_i^{-1}(\underline W)$, where $f_i=i_Ai^*$, so that successful sufficient conditions might also characterize the usual weak equivalences. The comparison diagram from pp.96--98 is rebuilt using relative asphericity of $A$, contractibility of the categories $i(a)$, preservation of weak equivalences by fibered categories under condition e), and a slice criterion f). The decisive computation identifies the slice over $(a,p)$ of $A_{/C}\to A_{//C}$ with the category of elements of a pullback of $i^*(\underline{Fl}(C))\to i^*(C)$; thus the problem becomes universal membership of this last map in $W_A$. This develops the earlier homotopy and internal-Hom machinery and supplies the lemma reformulated in section 62. The text records its own weakening on p.154: 0-admissibility b) is unnecessary for the intermediate comparison, since only $h_{(Cat)}\leq h_{\underline W}$ is needed. Two source slips require care. Diagram (11), p.153, ends in $A$, although the intended comparison and relative-asphericity argument require the projection $A\times C\to C$, as the earlier diagram on p.96 explicitly shows. The same paragraph changes the section reference to 57 while retaining pp.96--98; readers should follow those pages for the diagram. Neither is repaired in the readers. Page 155's literal $f_Y$ and lower mark by $Y$ are retained, and the visible label $p$ on the pullback square, omitted by the witness, is restored. At p.156 Grothendieck immediately withdraws the independent choice of the two arrow classes: the next section assumes $\underline W=\underline W'$ rather than claiming that this first axiomatization succeeded.

\subsection*{62. General localizers and the failure of the first characterization (pp.156--159)}
The new entry, visibly dated 6.3. rather than the witness's 6.4., retracts the attempt to choose two independent classes of arrows. With one saturated $\underline W$, the comparison diagram proves both $\underline W=f_i^{-1}(\underline W)$ and a canonical isomorphism of the induced $\overline f_i$ with the identity on $\underline W^{-1}(Cat)$; the middle arrow in the resulting chain is used through its inverse after localization. The hypotheses combine relative asphericity of $A$, contractibility of $i(a)$, the interval condition, the fibered-category condition, and the slice criterion developed in section 61. They require no restrictive 0-admissibility assumption, so even the class of all arrows satisfies them. This forces the main conceptual revision: the argument supplies a general modelizing theory rather than a unique characterization of ordinary weak equivalences. Grothendieck sketches coefficient-dependent, twisted and truncated cohomological variants, adopts $H_{\underline W}$ or $(Hot)_{\underline W}$ for the localized category, and asks whether ordinary weak equivalences form the smallest acceptable class, so that other examples would be localizations of the usual homotopy category. He also entertains finer notions, including simple homotopy types, as possibilities rather than results. These reflections motivate section 63's terminology and the later work with fundamental localizers. The marginal note on p.159 is consequential: it explicitly says that the impression is incorrect and suggests instead the biggest class under an extra requirement; the two words crossed by body typing remain doubtful in the checked text. The witness omits this qualification, as well as the inserted possibly before refinements, and therefore should not be used to turn the exploratory paragraph into a settled assertion. Three smaller source slips remain visible on p.157: the definition prints $i_Wi^*$ instead of the previously defined $i_Ai^*$, the three-map diagram is called (12) instead of (11), and saturation is called condition b) although it was a). The lower comparison row has primed objects, which the supplied witness row loses. The guide identifies these dependencies without changing the source formulas or replacing the author's subsequent change of mind with a modern theorem.

